% ======================================================================
\section{Geometry: rays, extinction and shielding columns}
\label{sec:geometry}
% ======================================================================

The MC field already contains the attenuation by dust (absorption and
scattering) and by the gas continua. What it does not contain is the
absorption in the \Htwo{} and CO lines, which the MC does not transport. That
absorption is applied through column densities integrated along the paths the
photons took: the straight line from each source for the direct beam, and the
path source $\to$ scattering point $\to$ cell for the scattered light. This
section describes how those paths and columns are built;
Sections~\ref{sec:h2model} and \ref{sec:photorates} apply them.

\subsection{Exact ray--cell traversal}
\label{sec:raycells}

Every ray in the bridge is traced by \rtn{rayCells}
(\src{pdr\_mod.f90:\allowbreak{}1710--1833}, \dlog{D-397}). MOCASSIN's cells are centred on
the axis values, so a cell's faces are the midpoints between neighbouring axis
values and the grid ends at the first and last values (the boundary cells are
half-cells). From a start point inside the grid and a unit vector $\hat u$,
the routine computes for each axis the distance along $\hat u$ to the next
face and repeatedly steps through the nearest one (a three-dimensional digital
differential analyser in the manner of \citealt{AmanatidesWoo1987}), returning
the list of cells crossed and the exact length $\ell$ of the ray in each, at a
cost proportional to the number of cells crossed. The first entry is the cell
the ray starts in; inactive cells (outside the nebula or in a cavity) are
returned with index zero and contribute nothing. The ray ends after a given
length or at the edge of the grid. For a \kw{symmetricXYZ} octant the lower
faces are planes of symmetry: when \rtn{lgMirror} is set the ray is reflected
there ($u_q\to-u_q$) and continues, a segment that re-enters the cell it just
left being merged with it; otherwise the ray ends at the edge. This replaced
fixed-step marches that sampled each cell an integer number of times and
missed the part of the neighbouring cell nearest the target, which had made
the source-ward $N(\mathrm{H_2})$ through \texttt{hiipdr}'s transition 0.46--0.75
of its exact value; with the exact paths that transition moved 0.25\,pc
towards the star (\dlog{D-397}, \dlog{H-025}, \dlog{L-45}).

\subsection{Source-ward extinction and molecular columns}
\label{sec:avradial}

\rtn{computeAVradial} (\src{pdr\_mod.f90:\allowbreak{}1859--2015}), called from
\rtn{computeAVdust} (\src{:\allowbreak{}2410--2523}) for point sources, traces for every
active cell and every source $s$ a ray from the cell centre towards the
source (no mirror) and accumulates
\begin{equation}
  \tau_{V,s} = \sum_{\mathrm{cells}} (\kappa^{\mathrm{abs}}_V+\kappa^{\mathrm{sca}}_V)\,\ell,
  \qquad
  N_{X,s} = \sum_{\mathrm{cells}\neq\mathrm{target}} n_{\mathrm H}\, x_X\,\ell,
  \quad X\in\{\mathrm{H_2},\,\mathrm C,\,\mathrm{CO}\},
\end{equation}
with the abundances left by the previous chemistry. $\tau_{V,s}$ runs to the
cell centre, the target's half-chord included. The target cell is excluded
from $N_{X,s}$ (\dlog{D-088}): including its near half made its \Htwo{}
shield itself, a positive feedback (more \Htwo, less dissociation, more \Htwo)
that, with MOCASSIN's large cells, drove the surface cell to complete
self-shielding. The column is therefore the one to the cell's near face. It
misses the cell's own contribution, an error that shrinks with the cell size
instead of feeding back: on \texttt{hiipdr}'s $31^3$ grid a cell is 0.648\,pc,
0.104\,mag of $\AV$ (\dlog{H-026}) and $2\times10^{20}$ H nuclei\,cm$^{-2}$, so
half a cell can hold $5\times10^{19}$\,cm$^{-2}$ of \Htwo{} and, at a CO
abundance of $10^{-4}$, $10^{16}$\,cm$^{-2}$ of CO, columns at which both
self-shield strongly (the CO factor of \citealt{vanDishoeckBlack1988} is
$\approx0.05$ at $N(\mathrm{CO})=10^{16}$\,cm$^{-2}$). The level model averages the transmission through the cell's own
column instead (Section~\ref{sec:ownavg}); the CO photorate and the
chemistry's own \Htwo{} diagnostic do not: their column lags the column to
the cell centre by half a cell (\dlog{D-088}), so they are under-shielded in
the cells of their own transitions on coarse grids, which places those
transitions deeper by up to about half a cell. The length of the source-ward ray inside
the target cell, from its near face to its centre, is kept separately as
$\ell_{\mathrm{own}}$ for the scattering sites (\dlog{D-398}).

The sources are combined with the geometric weights $f_s = r_s^{-2}$
(dilution only, \src{:\allowbreak{}1986}):
\begin{align}
  \AV &= -\frac{1.0857}{\RFUV}\,\ln\!\left[\frac{\sum_s f_s\, e^{-\RFUV\tau_{V,s}}}{\sum_s f_s}\right],
  &
  N_X &= \frac{\sum_s f_s N_{X,s}}{\sum_s f_s},
  &
  \ell_{\mathrm{own}} &= \frac{\sum_s f_s\,\ell_{\mathrm{own},s}}{\sum_s f_s}
  \label{eq:avradial}
\end{align}
(\src{:\allowbreak{}1996--2009}), stored as \rtn{AVeff}, \rtn{pdrColH2}, \rtn{pdrColC},
\rtn{pdrColCO} and \rtn{difSrc(:,6)}. For one source these are the plain
column and $\AV$, and all the fixtures have one source. For several, the $\AV$
is the one that reproduces the summed transmitted FUV of sources of
\emph{equal} luminosity, and the columns are means weighted by geometry alone
(neither the luminosities nor the attenuation of the sources enter $f_s$),
applied to the summed direct field through shielding functions that are not
linear in the column; the error of this is not measured. An outward ray along
the same direction (with mirrors) gives the dust column to the edge,
\rtn{AVback}.

$\AV$ is a coordinate and a diagnostic in the 3D build: the photorates do not
use it (Section~\ref{sec:rays}). It is a geometric dust column, not an
inversion of the MC field, which would under-read it by a factor 1.8 because
the field carries scattered light (\dlog{D-086}, \dlog{D-087}); TORUS-3DPDR also
derives $\AV$ from the column density along rays \citep[their Eq.~2]{Bisbas2015}.
It is measured from the source and so includes the dust of the ionised gas:
\texttt{hiipdr}'s ionisation front sits at $\AV=0.48$ (\dlog{D-260}). A 1D PDR
model, whose $\AV$ starts at its illuminated face, has to be compared at
$\AV$ measured from the front.

Under plane-parallel illumination \rtn{computeAVdust} integrates $\AV$ and the
columns along the beam direction itself (\src{:\allowbreak{}2432--2521}); there the
columns run to the cell centre, the cell's near half included
(\src{:\allowbreak{}2492--2507}), unlike the point-source branch. The level model does
not run in that geometry.

\subsection{Snapshot for the diffuse walk}

\rtn{snapDiffuseSources} (\src{pdr\_mod.f90:\allowbreak{}2082--2121}) is called by
\rtn{computeAVdust} before any cell is solved and stores per cell
(\rtn{difSrc}):
\begin{equation}
  n(\mathrm{H_2}),\ n(\mathrm C),\ n(\mathrm{CO})\ [\mathrm{cm^{-3}}];\qquad
  \kappa_{\mathrm{ext}} = \RFUV\,\kappa_V\ [\mathrm{cm^{-1}}];\qquad
  \epsilon = \sum_{k\in\mathrm{FUV}} \kappa^{\mathrm{sca}}_k\,\Jste_k\ [\mathrm{erg\,s^{-1}\,cm^{-3}}],
\end{equation}
the number densities that shield, a grey FUV extinction coefficient, and the
FUV power scattered per unit volume in the band (\src{:\allowbreak{}2099--2117}). For
isotropic scattering $\epsilon/4\pi$ is the band-integrated scattering
emissivity. It is computed from the full local MC field, direct and diffuse,
so it contains every order of scattering.

\subsection{Scattering sites of the diffuse field}
\label{sec:sites}

The diffuse FUV field is produced by scattering inside the cloud. A photon
scattered at a point $P$ and arriving at the cell $C$ has crossed the column
$N(\star\to P)+N(P\to C)$, not the column from the star to $C$
(\dlog{D-347}, \dlog{D-348}). \rtn{traceDiffuseSites}
(\src{pdr\_mod.f90:\allowbreak{}2136--2276}) represents the diffuse field arriving at the
target cell by a set of \emph{scattering sites}, built as the discretised
formal solution of the transfer equation for the intensity arriving at $C$
from each of a set of directions $\hat k$,
\begin{equation}
  4\pi I_{\hat k} = \int_0^\infty \epsilon(s)\, e^{-\tau(s)}\,\dd s ,
  \qquad \tau(s)=\int_0^s \kappa_{\mathrm{ext}}\,\dd s',
  \label{eq:formal}
\end{equation}
with $s$ the distance from $C$ along $\hat k$ (the direction the photons come
from):
\begin{enumerate}
\item \textbf{Directions.} The twelve HEALPix $N_{\mathrm{side}}=1$ pixel
  centres \citep{Gorski2005}, which divide the sphere into twelve equal solid
  angles of $\pi/3$\,sr, built directly in RING order: rays 0--3 at
  $\mu=+2/3$, $\phi=(k+\frac12)\pi/2$; rays 4--7 at $\mu=0$, $\phi=k\pi/2$;
  rays 8--11 at $\mu=-2/3$, $\phi=(k+\frac12)\pi/2$ (\src{:\allowbreak{}2152--2165}),
  with $\mu$ measured along the grid's $z$ axis. The directions are fixed in
  the grid frame.
\item \textbf{Walk.} Each ray is traced with \rtn{rayCells} from the target's
  centre, with mirror reflection at the octant's lower faces, to the edge of
  the model. Every cell crossed that scatters ($\epsilon>0$, $\kappa_{\mathrm{ext}}>0$)
  is a site, the target cell itself included (with the half-chord the ray
  crosses); the walk is stored (\rtn{wkCell},
  \rtn{wkLen}, \rtn{siteRay}, \rtn{siteIst}) for the level model.
\item \textbf{Weight.} With $\tau$ the grey FUV optical depth between the
  target's centre and the entry of the site cell and $\Delta\tau=\kappa_{\mathrm{ext}}\ell$
  across it,
  \begin{equation}
    w_{\mathrm{site}} = \epsilon\;\frac{1-e^{-\Delta\tau}}{\kappa_{\mathrm{ext}}}\;e^{-\tau}
    \qquad\bigl(\tfrac{1-e^{-\Delta\tau}}{\kappa_{\mathrm{ext}}}\to\ell(1-\tfrac12\Delta\tau)\ \text{for}\ \Delta\tau\le10^{-4}\bigr)
    \label{eq:siteweight}
  \end{equation}
  (\src{:\allowbreak{}2207--2217}): Equation~(\ref{eq:formal}) integrated over the
  site's segment with $\epsilon$ and $\kappa_{\mathrm{ext}}$ constant in the cell.
  Since the directions subtend equal solid angles,
  $4\pi J_{\mathrm{formal}}=\frac1{12}\sum_{\mathrm{sites}}w_{\mathrm{site}}$.
\item \textbf{Two-leg column.} For each species $X$ (\src{:\allowbreak{}2219--2239}),
  \begin{equation}
    N_{X,\mathrm{site}} = N_X^{\star}(P) + n_X(P)\,\ell_{\mathrm{own}}(P)
       + \sum_{\mathrm{cells\ between}} n_X\,\ell + \tfrac12\,n_X(P)\,\ell_P ,
    \label{eq:twoleg}
  \end{equation}
  the source-ward column of the site's cell to its near face, plus the site
  cell's own gas from its near face to its centre (\dlog{D-398}), plus the
  cells crossed between the target and the site, plus half of the site cell
  along the walk. The target cell is excluded from both legs, so a site in the
  target cell carries the direct beam's column.
\end{enumerate}

Only the \emph{relative} weights are used: the magnitude of the diffuse field
is the MC's $\chiF$ (Equation~\ref{eq:chisplit}), which the sites distribute
over directions and columns, each site becoming a ray of the chemistry with
field $\chiF\,w_{\mathrm{site}}/\sum w$ (Section~\ref{sec:rays}). The
distribution is the formal solution only under the assumptions built into
Equation~(\ref{eq:siteweight}): isotropic scattering, whereas the grains of
the fixtures have $g=0.61$ and 0.29 (Table~\ref{tab:dustfix}), so the real
diffuse field arrives preferentially from the source side; a grey extinction
equal to the source-weighted band mean, whereas the window in which \Htwo{} and
CO dissociate has a larger extinction and the local field at depth a smaller
one (Section~\ref{sec:dust}); an emissivity made by dust scattering alone,
whereas $\chiF$ also holds gas emission; uniform emissivity across a cell; and
the twelve-direction quadrature. As a check of these assumptions, the formal
solution's own magnitude, $4\pi J_{\mathrm{formal}}/(c\,\uD)$, recovers 52--130\%
of the MC's $\chiF$, falling with depth, for reasons not yet identified (gas
two-photon emission, the phase function and the quadrature were named as
candidates; the source-weighted grey extinction, which overstates the
extinction at depth, is another; \dlog{L-4}). The field arrives markedly
anisotropically: the brightest of the twelve directions carries 1.9--3.1 times
the mean (\dlog{D-348}). What the assumptions affect is the distribution of
the diffuse field over the sites' columns, not its magnitude, and that
distribution acts on the photorates only through the line shielding: at fixed
abundances the two-leg columns took back 5--32\% of the increase of the \Htwo{}
photodissociation rate that the diffuse field's magnitude produced
(\dlog{D-348}).

\subsection{Angular quadrature}
\label{sec:healpix}

The twelve directions of the sites and of the line rays are the coarsest
HEALPix tessellation. Their moments are not isotropic: the mean fourth power
of a direction's component is 0.132 along the ring axis ($z$) against 0.218
across it, where an isotropic set gives 0.2 (\dlog{D-410}). The result is a
pattern fixed in the grid frame: of the cells that the octant's symmetry makes
equivalent, the one nearest the plane $z=0$ holds about 11\% more \Htwo{}
before and across the H/\Htwo{} transition, on both fixtures, independently of
the packet count (\dlog{L-53}). Rotating the set moves the excess to the new
ring axis, and 48 directions ($N_{\mathrm{side}}=2$, fourth moments 0.188 and
0.201) remove it at 1.48 times the run time, moving every integrated quantity
by less than 1\%; the twelve were kept because no line map carries the pattern
(\dlog{D-410}). 3D-PDR and TORUS-3DPDR also use HEALPix rays, twelve
($\ell=0$) in TORUS-3DPDR's production runs \citep{Bisbas2012,Bisbas2015}.

\subsection{Line-escape rays}
\label{sec:lineRays}

\rtn{traceLineRays} (\src{pdr\_mod.f90:\allowbreak{}2327--2376}, \dlog{D-377}) traces the
same twelve directions from the target cell to the edge of the model, with
mirrors at the lower faces, skipping inactive cells, and returns for each ray
the cells crossed and the path length in each; the first entry is the cell
itself with the half of it the ray crosses. These rays give the escape
probabilities of the coolant lines (Section~\ref{sec:thermal}) and of the
\Htwo{} lines (Section~\ref{sec:escape}), each crossed cell contributing with
its own level populations. Before \dlog{D-377} the 3D flavour's line escape
was a slab with four of the twelve rays shut, which capped the escape
probability of an optically thin line at 2/3 (\dlog{L-27}).

\subsection{Comparison with 3D-PDR and TORUS-3DPDR}

In 3D-PDR \citep{Bisbas2012} the FUV field is a boundary condition: at every
point the field incident on the cloud surface is attenuated along HEALPix rays
traced to the boundary, $\chi(p)=N_\ell^{-1}\sum_{q}\chi_0(q)\,e^{-\tau_{\mathrm{UV}}\AV(q)}$
with $\tau_{\mathrm{UV}}=3.02$ in the benchmark comparisons
\citep[their Eq.~1]{Bisbas2015}, and the shielding columns are integrated
along the same rays. For a source inside the cloud these rays are the wrong
paths: a ray pointing away from the source would deliver unshielded radiation
from outside a nebula that has no outside source (\dlog{D-342}). TORUS-3DPDR
\citep{Bisbas2015} takes the FUV field of every cell from its MC path-length
estimator (912--2400\,\AA, stored as two vectors for the positive and negative
directions), derives $\AV$ and the column densities along HEALPix rays from
the density field, and leaves the dust's contribution to the diffuse field to
later work (their Sect.~3.1). The present coupling differs from both in that the dust
attenuation and scattering are the MC's, the field is split into its direct
and scattered parts, and the line-shielding columns follow the photons'
paths: source-ward for the direct beam and through the scattering sites for
the diffuse field.
